A positive integer $m$ is \emph{powerful} (a.k.a. \emph{squarefull} or \emph{$2$-full}) if for every prime $p\mid m$ we have $p^2\mid m$.

The prompt asks (paraphrasing) whether every pair of consecutive powerful integers $(n,n+1)$ must arise from a Pell-type construction, and then says ``in other words'' whether one of $n,n+1$ must be a square.

\textbf{Ambiguity / minimal correction.} ``Comes from a Pell equation'' is not a mathematically precise property unless one specifies what counts as ``coming from.'' The only precise embedded claim is:

\medskip
\noindent\textbf{(\*)} \emph{If $n$ and $n+1$ are powerful, then at least one of $n$ or $n+1$ is a perfect square.}

I will disprove (\*) by an explicit counterexample.
